\subsection{Rejection Masking Analysis}
\label{sec:masking_analysis}

A critical implementation detail in non-compact speculative decoding is whether rejected positions are zeroed in the attention mask before subsequent verification rounds.
We discovered a previously unreported interaction between masking and recurrent drafters.

\paragraph{Masking is necessary for correctness.}
We compare verifier outputs with and without rejection masking against gold autoregressive decoding.
With masking, the verifier's logit distribution differs from gold AR by only 0.02 (mean absolute diff); without masking, this divergence grows to 4.10.
The unmasked verifier attends to stale key-value entries from rejected tokens, producing incorrect distributions.

\paragraph{Recurrent and transformer drafters show opposite masking effects.}
Table~\ref{tab:mask_cross} reveals a striking asymmetry: the guided Mamba2 drafter achieves $\bar{n}_a\!=\!3.17$ without masking but only 2.33 with masking, while the LLaMA-1B transformer drafter shows the \emph{opposite} pattern (2.13 without, 2.73 with).

\begin{table}[t]
\centering
\caption{Cross-verification of rejection masking effects on acceptance rate ($\bar{n}_a$) for SSM vs.\ transformer drafters on UltraChat. SSM drafters show inflated acceptance without masking; transformer drafters show the opposite.}
\label{tab:mask_cross}
\vspace{0.5em}
\begin{tabular}{@{}lccc@{}}
\toprule
\textbf{Drafter} & \textbf{Masked} & \textbf{Unmasked} & \textbf{$\Delta$} \\
\midrule
Guided Mamba2-65M & $\sim$0.80 & 2.91 & +2.11 $\uparrow$ \\
LLaMA-3.2-1B     & 2.73     & 2.13 & $-$0.60 $\downarrow$ \\
\bottomrule
\end{tabular}
% TODO: Run full masked evaluation for Guided Mamba2 (currently only 16 samples)
% TODO: Add LLaMA-3B drafter row
\end{table}


This asymmetry arises because SSM drafters share no attention mechanism with the verifier.
Without masking, both the verifier and drafter process overlapping stale context through their respective state representations, creating artificial agreement.
Transformer drafters, maintaining their own KV cache, diverge from the verifier's corrupted state and thus show lower acceptance without masking.

\paragraph{The compact re-prefill bug.}
We further identified a critical issue where rejected draft tokens were incorrectly fed through the Mamba2 recurrence during re-prefill, contaminating the drafter's hidden state.
Fixing this bug improved masked acceptance from $\sim$0.9 to $\sim$2.2, a $2.3\times$ improvement.
All results in this paper are reported after this correction.

\subsection{Guidance Layer Ablation}
\label{sec:layer_ablation}

We ablate which verifier layers contribute most to acceptance in Table~\ref{tab:layer_ablation}.
We evaluate nine configurations with LLaMA-8B as verifier, varying the number and position of guidance extraction layers.

\begin{table}[t]
\centering
\caption{Guidance layer ablation (masked, greedy, LLaMA-8B verifier, Mamba2-27M drafter). All configurations with backbone finetuning show similar acceptance ($\sim$1.5\% spread), while frozen guidance underperforms.}
\label{tab:layer_ablation}
\vspace{0.5em}
\begin{tabular}{@{}llc@{}}
\toprule
\textbf{Configuration} & \textbf{v\_layers} & \textbf{Mean $\bar{n}_a$} \\
\midrule
KD $\to$ guided (default)     & [5, 16, 29] & \textbf{2.33} \\
Pair [5, 29]                   & [5, 29]     & 2.32 \\
z-branch [5, 16, 29]          & [5, 16, 29] & 2.30 \\
Single [29]                   & [29]        & 2.30 \\
Pretrain $\to$ guided + ft    & [5, 16, 29] & 2.30 \\
Pair [16, 29]                  & [16, 29]    & 2.30 \\
Single [16]                   & [16]        & 2.29 \\
\midrule
Pretrain $\to$ guided (frozen) & [5, 16, 29] & 1.98 \\
Single [5]                    & [5]         & 1.90 \\
\bottomrule
\end{tabular}
\end{table}

\paragraph{Layer selection barely matters when the backbone is finetuned.}
The top seven configurations span only 1.5\% in acceptance (2.29 to 2.33), indicating that the guidance signal is robust to layer choice when the backbone adapts alongside the projections.
Even a single high layer (layer 29, acceptance 2.30) nearly matches the full three-layer configuration (2.33).

\paragraph{Backbone finetuning is critical.}
Frozen guidance with three layers achieves only 1.98, while finetuned guidance reaches 2.30, representing $3\times$ the benefit of guidance projection alone.
This suggests that the backbone must adapt its representations to exploit guidance effectively.

\paragraph{KD is not mandatory.}
Starting from a pretrained (non-KD) backbone with finetuning achieves 2.30, nearly matching the KD-initialized variant (2.33).
This simplifies the training pipeline for practitioners who may not have access to a distillation setup.

\paragraph{z-branch gating adds no benefit.}
Injecting guidance into the gating pathway (z-branch) alongside the x-branch yields 2.30, no better than x-branch alone (2.33).
The state-update pathway is sufficient for guidance injection.

\subsection{Activation Replay Latency}
\label{sec:replay_results}

Table~\ref{tab:replay} compares activation replay against full re-prefill at various context lengths on CPU (Intel Xeon 8562Y+).

\begin{table}[t]
\centering
\caption{Activation replay vs.\ full re-prefill latency on CPU for Mamba2-27M. Replay is $O(K)$ in accepted tokens while re-prefill is $O(T)$ in context length.}
\label{tab:replay}
\vspace{0.5em}
\begin{tabular}{@{}rccc@{}}
\toprule
\textbf{Context length} & \textbf{Re-prefill (ms)} & \textbf{Replay (ms)} & \textbf{Speedup} \\
\midrule
32    & 299   & 47 & $6.4\times$ \\
128   & 1,019 & 37 & $27.7\times$ \\
512   & 3,644 & 36 & $100.9\times$ \\
1,024 & 8,036 & 40 & $201.0\times$ \\
\bottomrule
\end{tabular}
\end{table}

Re-prefill is sequential on CPU, costing $\sim$7\,ms per token $\times$ $T$ tokens.
Replay processes only the $\sim$5 accepted tokens, remaining roughly constant at $\sim$37\,ms regardless of context length.
At 1024 tokens, this yields a $201\times$ speedup, making activation replay essential for CPU-offloaded deployment.

On GPU, the benefit is less clear: the Mamba2-27M model is small enough that re-prefill completes in $\sim$16\,ms via parallel scan regardless of context, while sequential replay takes $\sim$40\,ms (5 kernel launches).
Activation replay is therefore a CPU-specific optimization in our setting.

\subsection{CPU Offloading Analysis}
\label{sec:cpu_results}

Table~\ref{tab:cpu_kernel} summarizes the CPU kernel optimization stack.

\begin{table}[t]
\centering
\caption{CPU inference optimization stack for Mamba2-27M (8-token draft, Intel Xeon 8562Y+ with AVX-512 VNNI).}
\label{tab:cpu_kernel}
\vspace{0.5em}
\begin{tabular}{@{}lrrc@{}}
\toprule
\textbf{Engine} & \textbf{8-tok draft} & \textbf{Throughput} & \textbf{vs.\ HF} \\
\midrule
HuggingFace PyTorch (8 threads)  & 79\,ms & 110\,tok/s & $1.0\times$ \\
AVX-512 FP32 SSM kernel          & 57\,ms & 123\,tok/s & $1.1\times$ \\
Fused BF16 (full fwd + LM head)  & 32\,ms & 340\,tok/s & $2.8\times$ \\
INT8 VNNI (fused)                 & 9\,ms  & 924\,tok/s & $\mathbf{7.5\times}$ \\
\bottomrule
\end{tabular}
\end{table}

\paragraph{Bottleneck analysis.}
Profiling reveals that the SSM selective scan accounts for only 10\% of per-layer latency; linear projections (in\_proj and out\_proj) dominate at 85\%.
This motivates our focus on fused BF16 and INT8 linear layers rather than SSM kernel optimization alone.
The AVX-512 SSM kernel achieves $4\times$ speedup over PyTorch for the scan operation, but the end-to-end gain is only $1.1\times$ because linear layers remain the bottleneck.

\paragraph{Pipeline overlap.}
With the INT8 kernel, 8-token draft latency (9\,ms at 16 threads) is strictly less than the LLaMA-8B verification time ($\sim$16\,ms on H100 NVL).
In the asynchronous pipeline, the GPU is the bottleneck: while it verifies one batch, the CPU has already completed the next draft.
With mean acceptance $\bar{n}_a\!=\!2.33$, the projected pipeline throughput is 202\,tok/s, or $2.7\times$ the 75.8\,tok/s autoregressive baseline.
This contrasts with GPU-only speculative decoding, which achieves $<1\times$ speedup because the drafter is too small to amortize the verification overhead on fast hardware.
